\def\biglen{20cm} % playing role of infinity (should be < .25\maxdimen)

\def\maxxy{3.3} % random points are in [-\maxxy,\maxxy]x[-\maxxy,\maxxy]

\begin{figure}[!ht]
    \centering
    \begin{subfigure}{.3\textwidth}
        \begin{tikzpicture}[scale=0.7]
        % \begin{scope}
            \def\radius{0.65}
            \def\pts{
                (-1.1483,1.74945),
                (0.8,-0.2),
                (1.5,-0.79076),
                (0.87991,-1.14056),
                (-0.9,-0.9)
            }
            % draw the points and their cells
            \xintForpair #1#2 in \pts \do{
                \edef\pta{#1,#2}
                \begin{scope}
                    \clip (\pta) circle (\radius);
                    \xintForpair \#3#4 in \pts \do{
                        \edef\ptb{#3,#4}
                        \ifx\pta\ptb\relax % check if (#1,#2) == (#3,#4) ?
                        \tikzset{myclip/.style={}}
                        \else
                        \tikzset{myclip/.style={clip}}
                        \fi
                        \path[myclip] (#3,#4) to[half plane] (#1,#2);
                    }
                    % \clip (-\maxxy,-\maxxy) rectangle (\maxxy,\maxxy); % last clip
                    % \pgfmathsetmacro{\randhue}{rnd}
                    % \definecolor{randcolor}{hsb}{\randhue,.5,1}
                    \fill[lightgray] (#1,#2) circle (\radius);
                    \node[dot, very thick] at (#1,#2) {}; % and draw the point
                \end{scope}
            }
            \coordinate (p1) at (-1.1483,1.74945);
            \coordinate (p3) at (1.5,-0.79076);
            \coordinate (p5) at (-0.9,-0.9);
            \coordinate (p4) at (0.87991,-1.14056);
            \coordinate (p2) at (0.8,-0.2);
        
            \fill[darkgray, opacity=0.3] (p2) -- (p3) -- (p4) -- cycle;
        
            \draw[thick] (p2) -- (p3) -- (p4) -- cycle;
        
            % \pgfresetboundingbox
            % \draw[opacity=0] (-\maxxy,-\maxxy) rectangle (\maxxy,\maxxy);
        % \end{scope}
        \end{tikzpicture}
        \subcaption*{$A_{1}$}
    \end{subfigure}\hfill
    \begin{subfigure}{.3\textwidth}
    \begin{tikzpicture}[scale=0.7]
    % \begin{scope}[xshift=7cm]
        \def\radius{0.89}
        \def\pts{
            (-1.1483,1.74945),
            (0.8,-0.2),
            (1.5,-0.79076),
            (0.87991,-1.14056),
            (-0.9,-0.9)
        }
        % draw the points and their cells
        \xintForpair #1#2 in \pts \do{
            \edef\pta{#1,#2}
            \begin{scope}
                \clip (\pta) circle (\radius);
                \xintForpair \#3#4 in \pts \do{
                    \edef\ptb{#3,#4}
                    \ifx\pta\ptb\relax % check if (#1,#2) == (#3,#4) ?
                    \tikzset{myclip/.style={}}
                    \else
                    \tikzset{myclip/.style={clip}}
                    \fi
                    \path[myclip] (#3,#4) to[half plane] (#1,#2);
                }
                % \clip (-\maxxy,-\maxxy) rectangle (\maxxy,\maxxy); % last clip
                % \pgfmathsetmacro{\randhue}{rnd}
                % \definecolor{randcolor}{hsb}{\randhue,.5,1}
                \fill[lightgray] (#1,#2) circle (\radius); % fill the cell with random color
                \node[dot, very thick] at (#1,#2) {}; % and draw the point
            \end{scope}
        }
        \coordinate (p1) at (-1.1483,1.74945);
        \coordinate (p3) at (1.5,-0.79076);
        \coordinate (p5) at (-0.9,-0.9);
        \coordinate (p4) at (0.87991,-1.14056);
        \coordinate (p2) at (0.8,-0.2);

        \fill[darkgray, opacity=0.3] (p2) -- (p3) -- (p4) -- cycle;

        \draw[thick] (p2) -- (p3) -- (p4) -- cycle;
        \draw[thick] (p4) -- (p5) -- (p2);

        % \pgfresetboundingbox
        % \draw[opacity=0] (-\maxxy,-\maxxy) rectangle (\maxxy,\maxxy);
    % \end{scope}
    \end{tikzpicture}
    \subcaption*{$A_{1.5}$}
    \end{subfigure}\hfill
    \begin{subfigure}{.3\textwidth}
    \begin{tikzpicture}[scale=0.7]
    % \begin{scope}[xshift=14cm]
        \def\radius{1.4}
        \def\pts{
            (-1.1483,1.74945),
            (0.8,-0.2),
            (1.5,-0.79076),
            (0.87991,-1.14056),
            (-0.9,-0.9)
        }
        % draw the points and their cells
        \xintForpair #1#2 in \pts \do{
            \edef\pta{#1,#2}
            \begin{scope}
                \clip (\pta) circle (\radius);
                \xintForpair \#3#4 in \pts \do{
                    \edef\ptb{#3,#4}
                    \ifx\pta\ptb\relax % check if (#1,#2) == (#3,#4) ?
                    \tikzset{myclip/.style={}}
                    \else
                    \tikzset{myclip/.style={clip}}
                    \fi
                    \path[myclip] (#3,#4) to[half plane] (#1,#2);
                }
                % \clip (-\maxxy,-\maxxy) rectangle (\maxxy,\maxxy); % last clip
                % \pgfmathsetmacro{\randhue}{rnd}
                % \definecolor{randcolor}{hsb}{\randhue,.5,1}
                \fill[lightgray] (#1,#2) circle (\radius); % fill the cell with random color
                \node[dot, very thick] at (#1,#2) {}; % and draw the point
            \end{scope}
        }
        \coordinate (p1) at (-1.1483,1.74945);
        \coordinate (p3) at (1.5,-0.79076);
        \coordinate (p5) at (-0.9,-0.9);
        \coordinate (p4) at (0.87991,-1.14056);
        \coordinate (p2) at (0.8,-0.2);

        \fill[darkgray, opacity=0.3] (p2) -- (p3) -- (p4) -- cycle;
        \fill[darkgray, opacity=0.3] (p2) -- (p4) -- (p5) -- cycle;
        \fill[darkgray, opacity=0.3] (p1) -- (p2) -- (p5) -- cycle;

        \draw[thick] (p2) -- (p3) -- (p4) -- cycle;
        \draw[thick] (p4) -- (p5) -- (p2);
        \draw[thick] (p1) -- (p2) -- (p5) -- cycle;

        % \pgfresetboundingbox
        % \draw[opacity=0] (-\maxxy,-\maxxy) rectangle (\maxxy,\maxxy);
    % \end{scope}
    \end{tikzpicture}
    \subcaption*{$A_{2}$}
    \end{subfigure}

    % \end{tikzpicture}
    \caption{\centering Esempio di complesso Alpha con decomposizione di Voronoi}
\end{figure}
